% translation-source: docs/user/latex/chapters/18-masques-retouche.tex; sha256: df922ca18b03a116c0eeb8f78f6170e72f0fc7864359dd88335cf82c0fc6bfab
\chapter{Masks and local editing}
The \emph{Masks} tab, at the top of the right panel, applies settings to one part of the image
only. A mask is a zone that carries its own values: they add to the global settings in that zone.
An exposure of +1 in a mask acts as a global exposure of +1, but only there. Nothing is written in
the photo: masks are saved in its sidecar and can be undone like the other steps.

Masks have neither order nor blend mode: where several overlap, their values add up.

\section{Creating a mask}
\begin{itemize}
\item \emph{Brush}: paint on the image. \emph{Size}, \emph{Feather} and \emph{Flow} set the next
stroke; each stroke is one history step.
\item \emph{Linear gradient}: the effect is whole before the first handle and none after the second.
\item \emph{Radial gradient}: move the centre and both edges on the image. \emph{Feather} sets the
transition, \emph{Angle} turns the ellipse.
\item \emph{Range}: selects tones of the image, between \emph{Dark tones} and \emph{Light tones},
as they are before any edit.
\end{itemize}
\emph{Show the zone} shows where the mask acts: in colour inside, in grey outside. This display is
neither saved nor exported. \emph{Invert} turns a shape over.

\section{Composing a zone}
\emph{Add to the zone} combines another shape with the first: \emph{Add}, \emph{Subtract} or
\emph{Intersect}. A range intersected with a gradient gives, for example, the shadows of the lower
part of the image alone.

\section{Adjusting a mask}
Each mask offers the light, colour, effects and detail settings, and an \emph{Opacity} that weighs
them all. The box in the list turns it off without losing it, and the two buttons beside it
duplicate or remove it. A mask follows the image's geometry: it stays on the same object after a
rotation or a crop.

Masks belong to their photo: pasting settings or applying them to a selection does not pass them
on. \emph{Reset settings} removes them too.

\section{Limits}
\begin{itemize}
\item A mask's luminance noise is processed after the profile, even in signed development.
\item Changing a mask is slower when there are many: beyond about fifteen, the gesture exceeds a
quarter of a second. Moving another slider does not draw them again.
\end{itemize}
